\section{Examples}\label{sec:examples}

The preceding criterion has elementary infinite-base examples.  They are
included only to illustrate the construction; no claim about a new arithmetic
system is intended.

\begin{example}[Natural numbers]
Let $D$ have carrier $\mathbb N$ with the usual addition and multiplication and
with division undefined when the denominator is zero.  In the free completion,
$a/0$ is therefore a non-base normal term rather than a natural number.
Addition satisfies
\[
 n+0=n\qquad(n\in\mathbb N).
\]
Taking the undefined term $0/0$ as $c_0$ and the context $x\mapsto x+0$ in
\cref{thm:base-fixing-properness} shows that
$\FCcomp{\Free D}$ is proper.
\end{example}

\begin{example}[Integers]
The same argument applies to $\mathbb Z$ with the usual ring operations and
partial division.  Again $x\mapsto x+0$ fixes the whole base, so the
finite-complement uniform completion is proper.
\end{example}

These examples should not be confused with totalized fields, meadows, common
meadows, or wheels, which impose algebraic conventions for singular division
\cite{BergstraHirshfeldTucker2009,BergstraPonseCommonMeadows,Carlstrom2004}.
Here no numerical value is assigned to $a/0$; the undefined application is
retained as a formal element of the free completion.
